\section{Revenue Model}

\label{sec:revenue}
% ============================================================================

Hanzo Platform generates revenue from five sources:

\begin{table}[H]
\centering
\small
\begin{tabular}{lrp{5cm}}
\toprule
\textbf{Source} & \textbf{Share} & \textbf{Description} \\
\midrule
Inference fees & 45\% & Per-token charges for Zen model inference \\
Training/fine-tuning & 20\% & GPU-hour charges for model customization \\
Platform subscriptions & 15\% & Tiered access to IAM, Base, Insights, etc. \\
Marketplace commissions & 10\% & Fees on third-party model listings \\
Enterprise contracts & 10\% & Custom SLAs, dedicated compute, on-prem \\
\bottomrule
\end{tabular}
\caption{Revenue composition (projected steady-state).}
\label{tab:revenue}
\end{table}

\subsection{Demand Projections}

Based on 14 months of production data (47M inference requests, 2,800+
applications), we model compute demand growth as:
\begin{equation}
D(t) = D_0 \cdot e^{\gamma t},
\end{equation}
where $D_0$ is baseline daily demand and $\gamma \approx 0.15$ (monthly growth
rate observed in production). At this rate, annual compute demand doubles
approximately every 5 months, consistent with industry-wide AI adoption
trends~\cite{mckinseyai2024}.

\subsection{Token Demand}

Organic token demand arises from three sinks:

\begin{enumerate}[leftmargin=1.4em]
    \item \textbf{Compute consumption}: Users must acquire HANZO to pay for
    inference and training. This creates continuous buy pressure proportional to
    platform usage.
    \item \textbf{Staking lock-up}: Node operators lock HANZO for the duration of
    their service, removing tokens from circulating supply.
    \item \textbf{Fee reduction lock-up}: API consumers hold HANZO in
    fee-reduction contracts, further reducing circulating supply.
\end{enumerate}

% ============================================================================
